\noindent For each $n\ge 1$, let $f(n)$ be the maximum size of a set $A\subseteq\{1,2,\dots,n\}$ with the property:

\begin{quote}
For all $a,b,c,d\in A$ with $a\le b\le c\le d$, if $abcd$ is a perfect square, then $ad=bc$.
\end{quote}

\noindent \emph{Ambiguity note.} Because $A$ is a set, elements of $A$ are distinct, but the statement does not explicitly say whether the quadruple $(a,b,c,d)$ is required to consist of four \emph{distinct} elements of $A$ or whether repetitions are allowed (e.g. $(a,b,b,c)$). The condition as written quantifies over $a,b,c,d\in A$ and allows equality in $a\le b\le c\le d$, so the literal reading allows repetitions. In what follows, I adopt this literal (stronger) reading; any set satisfying it also satisfies the ``distinct quadruples only'' variant.

\noindent The question asks for the asymptotic order of $f(n)$ (e.g. is $f(n)=o(n)$?).

\noindent 2)
